\section{Related work}\label{sec:related}

\paragraph{Statistical comparison of learning algorithms.}
Dietterich~\cite{dietterich1998} showed that the paired $t$-test on resampled hold-out scores has an inflated type~I error rate, because training sets overlap, and proposed the $5\times2$ cross-validation test. Nadeau and Bengio~\cite{nadeau2003} derived the corrected resampled $t$-test, which inflates the variance of repeated train/test differences by the factor $1/J + n_{\text{test}}/n_{\text{train}}$. Bengio and Grandvalet~\cite{bengio2004} proved that no unbiased estimator of the variance of $K$-fold cross-validation exists, which is why every practical procedure relies on a working model or a heuristic. Bouckaert and Frank~\cite{bouckaert2004} studied the replicability of these tests, Dem\v{s}ar~\cite{demsar2006} and Benavoli et al.~\cite{benavoli2017} treated comparisons over multiple data sets, frequentist and Bayesian. More recently, Bayle et al.~\cite{bayle2020} established central limit theorems for $K$-fold cross-validation under algorithmic stability, and Bates et al.~\cite{bates2024} clarified that cross-validation estimates the average error of the learning \emph{procedure} rather than that of the fitted model, and proposed nested cross-validation intervals. Bouthillier et al.~\cite{bouthillier2021} quantified the sources of variance (data split, initialisation, hyperparameter search) in machine learning benchmarks, and Varoquaux~\cite{varoquaux2018} showed how wide cross-validation error bars are in small samples. All these works assume independent observations.

\paragraph{Validation under spatial and temporal dependence.}
In ecology, remote sensing and the geosciences, a large literature shows that random cross-validation is optimistic for spatially structured data and recommends spatial blocking, buffering or nearest-neighbour distance matching~\cite{roberts2017,valavi2019,lerest2014,meyer2018,ploton2020,schratz2019,linnenbrink2024}. Wadoux et al.~\cite{wadoux2021} argued that, when the goal is the accuracy of a map over a fixed region, design-based probability sampling is preferable to any form of spatial cross-validation. In time series, Bergmeir and Ben\'{i}tez~\cite{bergmeir2012} compared blocked and out-of-sample schemes and Bergmeir et al.~\cite{bergmeir2018} showed when standard cross-validation remains valid for autoregressive models. These works are concerned with the \emph{point estimate} of predictive performance and its bias; the uncertainty of the comparison between two models, which is the object of this paper, is rarely addressed.

\paragraph{Comparing predictions of dependent data.}
Inference on the difference in accuracy between two predictions of dependent data is not new. The Diebold--Mariano (DM) test~\cite{diebold1995}, with the small-sample correction of Harvey et al.~\cite{harvey1997}, applies a HAC (Newey--West~\cite{newey1987}) variance estimator to the loss-differential series of two forecasts. Hering and Genton~\cite{hering2011} proposed the spatial prediction comparison test, which models the covariance of a spatially correlated loss-differential field with a parametric variogram, and Gilleland et al.~\cite{gilleland2018} showed that the size of such tests is sensitive to serial dependence and to the correlation between competing forecasts. Watson et al.~\cite{watson2024} studied prediction and model evaluation for space--time data, with emphasis on how the validation design must match the prediction task. HAC estimation under cross-sectional dependence goes back to Conley~\cite{conley1999} and, for panels, to Driscoll and Kraay~\cite{driscoll1998}; the poor size of HAC tests in small effective samples motivated the fixed-$b$ theory of Kiefer and Vogelsang~\cite{kiefer2005}. Block and stationary bootstraps~\cite{kunsch1989,politis1994,lahiri2003} provide resampling alternatives.

\paragraph{Positioning of this work.}
In view of this literature, our contribution is not the idea of accounting for dependence in the loss differential, which the DM and Hering--Genton tests already do for \emph{fixed} predictors. It lies in four more specific elements. (i)~We compare \emph{learning procedures} that are refitted on each training fold of a blocked cross-validation, which adds a training-induced, fold-level variance component, correlated across folds, that is absent from forecast or map comparisons; we make its estimation and the conditions for its validity explicit (Propositions~\ref{prop:pairedt} and~\ref{prop:bias}). (ii)~DA-P combines this component with the test-sampling dependence in a single variance estimate computed from one cross-validation run. (iii)~The dependence of the test-sampling fluctuations is modelled on the variograms of the \emph{fold-centred} loss differentials, so that it is not confounded with the training effects, and its adequacy is examined under several misspecifications. (iv)~The procedure returns diagnostics and a three-way decision that declares a comparison indeterminate when the data contain too few effectively independent regions; the threshold of this rule is calibrated and validated by simulation. We include a space--time version of the Hering--Genton test among the competing procedures in all our experiments.

\paragraph{Trustworthy machine learning and this special issue.}
Recent contributions in this journal have advocated statistical tools for the robustness of machine learning models~\cite{giudici2024}, practical safeguards against overfitting~\cite{chicco2026}, and effect-size reporting beyond $p$-values in large data sets~\cite{caponero2025}. The present paper complements them by addressing the validity of the comparison step itself when observations are dependent.
